\documentclass[11pt]{article}
\usepackage{nopageno}
\setcounter{secnumdepth}{1}
\usepackage[margin=1in]{geometry}
\begin{document}

\pagestyle{empty}

\noindent \textbf{Facilities, Equipment, and Other Resources:}

\begin{itemize}

\item \textbf{Heritage 3D Lab at Penn State:}
The Penn State Heritage 3D Lab, led by PI Napolitano, is a research facility dedicated to advancing the understanding and preservation of heritage structures through innovative engineering and computational methods. The lab brings together faculty, graduate and undergraduate students for research on structural analysis, heritage preservation, and disaster resilience. The Heritage 3D Lab houses field-deployable equipment supporting both research and teaching, including a Trimble X7 terrestrial laser scanner, Leica laser rangefinder, panoramic photosphere camera, various environmental sensors, and thermal imaging cameras. These tools enable on-site assessments of heritage structures and their environments, including the Montpelier field damage surveys that produced the 271-building dataset used in this proposal.

The lab also features an advanced workspace designed for integrated, team-based workshops. The ICon Lab features an immersive curved screen with ultra-high-definition projection capabilities, supporting 3D visualization and walkthroughs of architectural models. This space facilitates collaborative design, review, and training sessions for up to 30 people, and will host the hybrid community co-development workshops described in Subtask~2.1.

\item \textbf{Busse Lab at Penn State:}
The Busse Lab, led by Co-PI Busse, has four Dell Precision 3660 desktop computers dedicated to computational work, equipped with HEC-RAS 2D, EPANET, Python, MATLAB, and Solidworks. These resources directly support the mass balance modeling and hydrological analysis in Subtasks~1.2 and~2.4. The main lab space is equipped for water treatment and water system testing applications and is approved for BSL2-level experiments. Co-PI Busse also shares access to a closed-loop water channel built by Engineering Laboratory Design, Inc.

\item \textbf{Institute for Computational and Data Sciences -- Advanced Cyber Infrastructure:}
Research data will be hosted by Penn State's Institute for Computational and Data Sciences (ICDS) through its Roar supercomputer. Roar provides active storage (DDN 12KX40 and GS7K flash storage arrays) and near-line storage (Oracle FS1 flash appliance and SL8500 Tape Library), with daily backups. ICDS provides long-term archival services; research data will be archived for at least three years after the end of the award or after public release, whichever comes later.
Data Security: Roar requires strong password and two-factor authentication, with auditable access. Automated weekly vulnerability scans, data encryption in-flight and at rest, SSH-only firewalls, and least-privilege access controls protect research data. Roar storage is physically secured in Penn State's Tower Road Data Center (TRDC), with swipe-card access and continuous monitoring.

\item \textbf{Survey Research Center (SRC):}
The Survey Research Center at Penn State offers services to design and implement data collection methods, including internet surveys via Qualtrics or REDCap and support for focus groups and in-person fieldwork. SRC will support the structured community elicitation instruments used in the Subtask~2.1 workshops.

\item \textbf{Teleconference and Web Meeting Facility:}
Penn State provides teleconference and video capability free to the research team, supporting hybrid workshop delivery (Subtask~2.1), monthly PI coordination meetings, and webinar-based dissemination of research outcomes to SHPOs and other practitioners (Broader Impacts).

\item \textbf{Community, State, and Federal Agency Collaborators:}
The project leverages established relationships with local, state, and federal partners to ensure robust community engagement and alignment with national preservation and emergency management practice.
\begin{itemize}
    \item \textbf{Montpelier Alive (Katie Trautz, Executive Director):} Downtown merchant and community network in the state capital. Provides historic programmatic archives and local outreach capacity; community liaisons support building-owner participation in co-development workshops.
    \item \textbf{Foundation for a Resilient Montpelier (Jon Copans):} Local governance and recovery organization. Provides disaster recovery records and institutional frameworks; assists in translating risk outputs into municipal-scale policy proposals during co-production phases.
    \item \textbf{Vermont State Historic Preservation Officer (Laura Trieschmann):} Provides institutional context, historic preservation expertise, and access to historic district records for the Montpelier case study. Participates in the expert panel evaluating decision-support outputs.
    \item \textbf{U.S.\ Department of the Interior (Jennifer Wellock):} Co-author of the NPS Guidelines on Flood Adaptation for Historic Buildings. Provides expert oversight of building intervention typologies and facilitates nationwide dissemination of the transferable protocol to heritage-at-risk communities.
    \item \textbf{Association for Preservation Technology International -- Disaster Response Initiative (Joe Kallas):} APT DRI, which PI Napolitano co-directs, serves as the practitioner platform for field-testing and hosting the decision interface. APT DRI will host tabletop simulation exercises at its annual conference to evaluate practitioner decision-making using framework outputs.
\end{itemize}

\item \textbf{PSU Career Services:}
Graduate Career Services offers career counseling, mock interviewing, and professional development programs supporting the postdoctoral researcher funded by this project.

\item \textbf{PSU Graduate Writing Center:}
The Graduate Writing Center provides free one-to-one peer consultations and workshops for researchers at all levels, supporting the postdoctoral researcher's scholarly development and grant writing training.

\item \textbf{Penn State Center for the Integration of Research, Teaching and Learning (CIRTL):}
The CIRTL Network offers courses, workshops, and events on inclusive teaching, mentoring, and career development for postdoctoral researchers. Relevant courses include Research Mentor Training and Advancing Learning Through Evidence-Based STEM Teaching.

\item \textbf{PSU Scholarship and Research Integrity Education Program:}
The SARI program offers comprehensive, multilevel training in the responsible conduct of research (RCR). The postdoctoral researcher will complete SARI requirements covering research ethics, data management, and intellectual property.

\item \textbf{PSU Undergraduate Research and Fellowships Mentoring (URFM) Program:}
URFM provides resources on undergraduate research mentorship development and best practices, supporting the BESURE undergraduate researchers participating in this project.

\end{itemize}
\end{document}
